% Unit 086 English translation of chapter6.tex, lines 2604--2927.
% Source: Methods of Algebra, Volume 2, commit
% 9a5803ff77dd3257484cb177f851a73770a59dd3 (CC BY 4.0).
% stable-unit-id: o014.aljabr2.chapter6.nonabelian-cohomology
% Coverage: complete Section 6.13, ``Nonabelian Cohomology''; ends immediately
% before the Exercises environment at line 2928.

% segment-id: o014.li2.chapter6.nonabelian-cohomology.h001
\section{Nonabelian Cohomology}\label{sec:nonabelian-coh}
\hypertarget{o014.aljabr2.chapter6.nonabelian-cohomology}{ }
% segment-id: o014.li2.chapter6.nonabelian-cohomology.p001
Let $G$ be a group. The cohomology $\Hm^n(G,M)$ discussed previously always
requires its ``coefficients'' $M$ to have a module structure, or at least to
be a $\Z$-module, that is, an abelian group. This section studies the case in
which the coefficients form a nonabelian group. From now on we use the symbol
$A$ in place of $M$.

% segment-id: o014.li2.chapter6.nonabelian-cohomology.d001
\begin{definition}\label{def:G-group}
	\index[sym1]{G-Grp@$G\dcate{Grp}$}
	\index{$G$-group}
	Let $A$ be a group, with its binary operation written multiplicatively
	and its identity element denoted by $1$. Suppose that $A$ is equipped with
	a left action of the group $G$, written
% segment-id: o014.li2.chapter6.nonabelian-cohomology.q001
	\[ G \times A \to A, \quad (g, a) \mapsto {}^g a . \]
	If the group action is compatible with the multiplicative structure of
	$A$,
% segment-id: o014.li2.chapter6.nonabelian-cohomology.q002
	\[ {}^g (a_1 a_2) = {}^g a_1 {}^g a_2, \quad {}^g 1 = 1, \]
	then $A$ is called a group with a $G$-action, or simply a
	\emph{$G$-group}. A homomorphism of $G$-groups is defined to be a group
	homomorphism that preserves the $G$-action. The resulting category is
	denoted by $G\dcate{Grp}$.
\end{definition}

% segment-id: o014.li2.chapter6.nonabelian-cohomology.p002
We will soon see that, for a $G$-group $A$, without additional structure or
commutativity only $\Hm^0$ and $\Hm^1$ can be defined in a reasonable and
useful way. In this setting $\Hm^1$ need not be a group; it is instead a
pointed set.

% segment-id: o014.li2.chapter6.nonabelian-cohomology.d002
\begin{definition}\label{def:pointed-set}
	\index{pointed set}
	\index[sym1]{Set-pt@$\cate{Set}_\bullet$, $G\dcate{Set}_\bullet$}
	A \emph{pointed set} is data $(X,x_0)$, where $X$ is a set and
	$x_0\in X$ is its basepoint. A morphism from $(X,x_0)$ to $(Y,y_0)$ is a
	map $f:X\to Y$ satisfying $f(x_0)=y_0$. These data form the category
	$\cate{Set}_\bullet$.

	Let $(X,x_0)$ be a pointed set. If a group $G$ acts on $X$ from the left,
	written $(g,x)\mapsto{}^gx$, and ${}^gx_0=x_0$ always holds, we say that
	$G$ acts on $(X,x_0)$. In that case we can define the pointed subset of
	invariants
% segment-id: o014.li2.chapter6.nonabelian-cohomology.q003
	\[ (X, x_0)^G := \left\{ x \in X: \forall g \in G, \; gx = x \right\}, \quad \text{basepoint} = x_0. \]
	All pointed sets with a $G$-action, together with morphisms compatible
	with that action, form the category denoted by $G\dcate{Set}_\bullet$.
\end{definition}

% segment-id: o014.li2.chapter6.nonabelian-cohomology.p003
When no confusion can arise, we often write an object $(X,x_0)$ of
$\cate{Set}_\bullet$ simply as $X$, and in $G\dcate{Set}_\bullet$ write
$(X,x_0)^G$ as $X^G$.

% segment-id: o014.li2.chapter6.nonabelian-cohomology.p004
Every group becomes a pointed set by choosing its identity element as the
basepoint. This gives the evident functor
$\cate{Grp}\to\cate{Set}_\bullet$; similarly, there is a functor
$G\dcate{Grp}\to G\dcate{Set}_\bullet$. Note that the invariant subset
$A^G$ of a $G$-group $A$ is a subgroup.

% segment-id: o014.li2.chapter6.nonabelian-cohomology.d003
\begin{definition}\label{def:pointed-exact}
	\index{exact sequence}
	Let
	$(X,x_0)\xrightarrow{f}(Y,y_0)\xrightarrow{g}(Z,z_0)$ be morphisms in
	$\cate{Set}_\bullet$ (or $G\dcate{Set}_\bullet$). This sequence of
	morphisms is called exact if $g^{-1}(z_0)=\Image(f)$. In this way, exact
	sequences in $\cate{Set}_\bullet$ (or $G\dcate{Set}_\bullet$) are
	defined.
\end{definition}

% segment-id: o014.li2.chapter6.nonabelian-cohomology.p005
For abelian groups this recovers the familiar notion of exactness. We can now
define $\Hm^0$ and $\Hm^1$ for every $G$-group $A$.

% segment-id: o014.li2.chapter6.nonabelian-cohomology.d004
\begin{definition}\label{def:nonabelian-H}
	\index{group cohomology!nonabelian case}
	Let $A$ be an object of $G\dcate{Set}_\bullet$. Define the pointed set
% segment-id: o014.li2.chapter6.nonabelian-cohomology.q004
	\[ \Hm^0(G, A) = Z^0(G, A) := A^G. \]

	If $A$ is a $G$-group, then $\Hm^0(G,A)$ is a group. In this case,
	further define
% segment-id: o014.li2.chapter6.nonabelian-cohomology.q005
	\begin{align*}
		Z^1(G, A) & := \left\{\begin{array}{r|l}
			c: G \to A & \forall g_1, g_2 \in G, \\
			& c(g_1 g_2) = c(g_1) \cdot {}^{g_1} c(g_2)
		\end{array}\right\}, \\
		\Hm^1(G, A) & := Z^1(G, A) / \sim ,
	\end{align*}
	where the equivalence relation $c\sim c'$ means that there is an
	$a\in A$ such that, for every $g\in G$,
% segment-id: o014.li2.chapter6.nonabelian-cohomology.q006
	\[ c'(g) = a^{-1} \cdot c(g) \cdot {}^g a . \]
	It is easy to see that this is indeed an equivalence relation---it comes
	from a right $A$-action---and $\Hm^1(G,A)$ naturally becomes a pointed
	set. Its basepoint is the equivalence class determined by the constant
	map $c(g)=1$.
\end{definition}

% segment-id: o014.li2.chapter6.nonabelian-cohomology.p006
Note that substituting $g_1=g_2=1_G$ into the condition defining
$Z^1(G,A)$ gives $c(1_G)=1$. Henceforth $[c]$ denotes the equivalence class
of $c\in Z^1(G,A)$.

% segment-id: o014.li2.chapter6.nonabelian-cohomology.r001
\begin{remark}[$\Hm^1$ as conjugacy classes of sections]\label{rem:Z1-vs-section}
	Form the semidirect product $A\rtimes G$ from the $G$-group $A$, and let
	$\pi$ denote the projection homomorphism $A\rtimes G\to G$. The elements
	of $Z^1(G,A)$ correspond bijectively to group homomorphisms
	$\sigma:G\to A\rtimes G$ satisfying $\pi\sigma=\identity_G$ (that is,
	``sections'' of $\pi$), via the map sending $c\in Z^1(G,A)$ to
	$\sigma(g)=(c(g),g)$. Moreover, $[c']=[c]$ if and only if there is an
	$a\in A$ such that $\sigma'=a^{-1}\sigma a$, where $A$ is embedded as a
	subgroup of $A\rtimes G$ and conjugation is performed pointwise.
\end{remark}

% segment-id: o014.li2.chapter6.nonabelian-cohomology.p007
If $A$ is an abelian group, these constructions reduce to the descriptions
of $\Hm^0$ and $\Hm^1$ using the standard cochain complex
(Proposition~\ref{prop:groupcoh-cochain}). In the nonabelian case, the
constructions remain functorial, as follows.
% segment-id: o014.li2.chapter6.nonabelian-cohomology.l001
\begin{itemize}
% segment-id: o014.li2.chapter6.nonabelian-cohomology.i001
	\item Every morphism $f:A_1\to A_2$ induces
% segment-id: o014.li2.chapter6.nonabelian-cohomology.q007
	\begin{gather*}
		\Hm^0(f): \Hm^0(G, A_1) \to \Hm^0(G, A_2), \quad
		\Hm^1(f): \Hm^1(G, A_1) \to \Hm^1(G, A_2);
	\end{gather*}
	the first map is evident, while the second is given by
	$\Hm^1(f)([c])=[f\circ c]$.

% segment-id: o014.li2.chapter6.nonabelian-cohomology.i002
	\item Let $\varphi:H\to G$ be a group homomorphism. Pull the $G$-action
	on the $G$-group $A$ back to $H$, giving the $H$-group $\varphi^*A$.
% segment-id: o014.li2.chapter6.nonabelian-cohomology.l002
	\begin{compactitem}
% segment-id: o014.li2.chapter6.nonabelian-cohomology.i003
		\item Define the canonical morphism
		$\Hm^0(G,A)\to\Hm^0(H,\varphi^*A)$ to be the evident inclusion
		$A^G\hookrightarrow A^{\varphi(H)}=(\varphi^*A)^H$.
% segment-id: o014.li2.chapter6.nonabelian-cohomology.i004
		\item Define the canonical morphism of pointed sets
% segment-id: o014.li2.chapter6.nonabelian-cohomology.q008
		\[ \Hm^1(G, A) \to \Hm^1(H, \varphi^* A). \]
		It is given by $[c]\mapsto[c\circ\varphi]$.
	\end{compactitem}
\end{itemize}

% segment-id: o014.li2.chapter6.nonabelian-cohomology.p008
These canonical maps have the various properties seen in
\S\ref{sec:change-of-group}. They can also be summarized as follows. Given
$\varphi:H\to G$, let $A_1$ be an object of $G\dcate{Set}_\bullet$ (or
$G\dcate{Grp}$), let $A_2$ be an object of $H\dcate{Set}_\bullet$ (or
$H\dcate{Grp}$), and let the morphism $f:A_1\to A_2$ be equivariant with
respect to $\varphi$, as in Definition~\ref{def:group-equivariance}. There
are corresponding maps $\Hm^i(f):\Hm^i(G,A_1)\to\Hm^i(H,A_2)$, and the
analogue of Proposition~\ref{prop:grp-equiv} takes the form:
% segment-id: o014.li2.chapter6.nonabelian-cohomology.l003
\begin{itemize}
% segment-id: o014.li2.chapter6.nonabelian-cohomology.i005
	\item $\Hm^0(f)$ is the restriction of $f$ to $A_1^G\to A_2^H$;
% segment-id: o014.li2.chapter6.nonabelian-cohomology.i006
	\item $\Hm^1(f)$ maps $[c]$ to $[f\circ c\circ\varphi]$.
\end{itemize}

% segment-id: o014.li2.chapter6.nonabelian-cohomology.c001
\begin{convention}
	Write $1$ for the zero object in $\cate{Set}_\bullet$ or
	$G\dcate{Set}_\bullet$ determined by the singleton set.
\end{convention}

% segment-id: o014.li2.chapter6.nonabelian-cohomology.p009
We now come to the main subject. Henceforth consider a short exact sequence
in $G\dcate{Set}_\bullet$ (Definition~\ref{def:pointed-exact})
% segment-id: o014.li2.chapter6.nonabelian-cohomology.q009
\[ 1 \to A \xrightarrow{u} B \xrightarrow{v} C \to 1, \]
where $A$ and $B$ are assumed to be $G$-groups and $u$ is a group
homomorphism. Exactness therefore implies that $u$ is injective and $v$ is
surjective. We further require:
% segment-id: o014.li2.chapter6.nonabelian-cohomology.q010
\begin{equation}\label{eqn:pointed-ses-condition}
	\begin{gathered}
		\forall b, b' \in B, \; \left[ v(b) = v(b') \iff \exists a \in A, \; b' = b u(a) \right], \\
		\text{in other words: $C$ is identified with $B/u(A)$, basepoint $= 1 \cdot u(A)$.}
	\end{gathered}
\end{equation}

% segment-id: o014.li2.chapter6.nonabelian-cohomology.p010
If $C$ is a group and $v$ is a group homomorphism, considering
$b^{-1}b'\in v^{-1}(1)=\Image(u)$ shows that
\eqref{eqn:pointed-ses-condition} holds automatically, and in this case
$u(A)\lhd B$. Many applications, however, do not fall into this case.

% segment-id: o014.li2.chapter6.nonabelian-cohomology.p011
Our present goal is to construct from these data an exact sequence that is
``as long as possible,'' together with the corresponding connecting
morphisms. Without further comment, we identify $A$ with its image as a
subgroup of $B$ and omit the symbol $u$.

% segment-id: o014.li2.chapter6.nonabelian-cohomology.p012
\textbf{The degree-zero case.}\;
First define the connecting morphism
$\delta^0:\Hm^0(G,C)\to\Hm^1(G,A)$. Let $c\in C^G$. Choose any
$b\in v^{-1}(c)$. Since $v$ preserves the $G$-action,
\eqref{eqn:pointed-ses-condition} shows that, for every $g\in G$, there is
a unique $a(g)\in A$ such that ${}^gb=ba(g)$. Hence $g\mapsto a(g)$ gives
an element $a$ of $Z^1(G,A)$, because
% segment-id: o014.li2.chapter6.nonabelian-cohomology.q011
% Editorial correction: the frozen Chinese source and Indonesian witness have
% a(g) in the final expression; the cocycle identity requires a(g_1).
\[ a(g_1 g_2) = b^{-1} \cdot {}^{g_1 g_2} b = b^{-1} \cdot {}^{g_1} b \cdot {}^{g_1} (b^{-1} \cdot {}^{g_2} b) = a(g_1) \cdot {}^{g_1} a(g_2). \]
If $b',b\in v^{-1}(c)$, there is an $\alpha\in A$ such that
$b'=b\alpha$, and hence
% segment-id: o014.li2.chapter6.nonabelian-cohomology.q012
\[ (b')^{-1} \cdot {}^g b' = \alpha^{-1} \cdot b^{-1} \cdot {}^g b \cdot {}^g \alpha. \]
This shows that the class $[a]\in\Hm^1(G,A)$ is determined solely by $c$.
If $c$ is the basepoint of $C$, we may take $b=1\in B$. We have therefore
obtained a morphism in $\cate{Set}_\bullet$,
$\delta^0:\Hm^0(G,C)\to\Hm^1(G,A)$.

% segment-id: o014.li2.chapter6.nonabelian-cohomology.p013
\textbf{The degree-zero case: a central subgroup.}\;
Suppose that $C$ is a group, $v$ is a group homomorphism, and $A$ is
contained in the center $Z_B$ of the group $B$; in particular, $A$ is
abelian. In this case $\Hm^0(G,C)=C^G$ and $\Hm^1(G,A)$ are both groups.
From the definition of $\delta^0$ and the condition $A\subset Z_B$, it is
easy to check that $\delta^0$ is a group homomorphism.

% segment-id: o014.li2.chapter6.nonabelian-cohomology.p014
\textbf{The degree-one case.}\;
Continue to suppose that $C$ is a group, $v$ is a group homomorphism, and
$A\subset Z_B$; then $\Hm^2(G,A)$ can be described using the standard
cochain complex. We define $\delta^1:\Hm^1(G,C)\to\Hm^2(G,A)$. Let
$c\in Z^1(G,C)$. For every $g\in G$, choose
$b(g)\in v^{-1}(c(g))$. Since \eqref{eqn:pointed-ses-condition} always
holds and $v$ preserves the $G$-action, there is a unique map $a:G^2\to A$
satisfying
% segment-id: o014.li2.chapter6.nonabelian-cohomology.q013
\[ b(g_1) \cdot {}^{g_1} b(g_2) = a(g_1, g_2) b(g_1 g_2). \]

% segment-id: o014.li2.chapter6.nonabelian-cohomology.p015
We now check that $a\in Z^2(G,A)$. This is equivalent to
% segment-id: o014.li2.chapter6.nonabelian-cohomology.q014
\[ a(g_1, g_2)^{-1} \cdot a(g_1, g_2 g_3) \cdot a(g_1 g_2, g_3)^{-1} \cdot {}^{g_1} a(g_2, g_3) = 1. \]
Expand the left-hand side as an expression in $B$:
% segment-id: o014.li2.chapter6.nonabelian-cohomology.q015
\begin{multline*}
	\underbracket{b(g_1 g_2) \cdot {}^{g_1} b(g_2)^{-1} \cdot b(g_1)^{-1}}_{a(g_1, g_2)^{-1}} \cdot \underbracket{b(g_1) \cdot {}^{g_1} b(g_2 g_3) b(g_1 g_2 g_3)^{-1}}_{a(g_1, g_2 g_3)} \\
	\underbracket{b(g_1 g_2 g_3) \cdot {}^{g_1 g_2} b(g_3)^{-1} \cdot b(g_1 g_2)^{-1}}_{a(g_1 g_2, g_3)^{-1}} \; \cdot \; {}^{g_1} a(g_2, g_3),
\end{multline*}
and then use $A\subset Z_B$ to simplify it to
% segment-id: o014.li2.chapter6.nonabelian-cohomology.q016
\begin{multline*}
	b(g_1 g_2) \cdot {}^{g_1} b(g_2)^{-1} \cdot {}^{g_1} b(g_2 g_3) \cdot {}^{g_1 g_2} b(g_3)^{-1}  b(g_1 g_2)^{-1} \cdot {}^{g_1} a(g_2, g_3) \\
	= b(g_1 g_2) \cdot {}^{g_1} b(g_2)^{-1} \cdot \underbracket{{}^{g_1} b(g_2) \cdot {}^{g_1 g_2} b(g_3) \cdot {}^{g_1} b(g_2 g_3)^{-1}}_{{}^{g_1} a(g_2, g_3)} \cdot {}^{g_1} b(g_2 g_3) \cdot {}^{g_1 g_2} b(g_3)^{-1}  b(g_1 g_2)^{-1} .
\end{multline*}
The last expression is visibly equal to $1$, so $a\in Z^2(G,A)$.

% segment-id: o014.li2.chapter6.nonabelian-cohomology.p016
If the choice of every $b(g)$ is changed and written
$b'(g)=\alpha(g)b(g)$, where $\alpha:G\to A$ is any map, then the
corresponding $a'(g_1,g_2)$ is
% segment-id: o014.li2.chapter6.nonabelian-cohomology.q017
\begin{align*}
	a'(g_1, g_2) & = \alpha(g_1) \cdot b(g_1) \cdot {}^{g_1} \alpha(g_2) \cdot {}^{g_1} b(g_2) b(g_1 g_2)^{-1} \alpha(g_1 g_2)^{-1} \\
	& = \alpha(g_1) \cdot {}^{g_1} \alpha(g_2) \cdot \alpha(g_1 g_2)^{-1} \cdot a(g_1, g_2) \quad (\because\; A \subset Z_B).
\end{align*}
Thus $a,a'\in Z^2(G,A)$ differ only by a factor belonging to $B^2(G,A)$.

% segment-id: o014.li2.chapter6.nonabelian-cohomology.p017
Finally, consider the case in which $c$ is replaced by
$c'(g)=\gamma^{-1}\cdot c(g)\cdot{}^g\gamma$, with $\gamma\in C$. Choose
any $\beta\in v^{-1}(\gamma)$. For the $b'$ corresponding to $c'$, we may
take $b'(g)=\beta^{-1}\cdot b(g)\cdot{}^g\beta$. The map $a':G^2\to A$
constructed from $b'$ is
% segment-id: o014.li2.chapter6.nonabelian-cohomology.q018
\begin{align*}
	a'(g_1, g_2) & = \left( \beta^{-1} \cdot b(g_1) \cdot {}^{g_1} \beta \right) \cdot {}^{g_1} \left( \beta^{-1} \cdot b(g_2) \cdot {}^{g_2} \beta \right) \cdot \left( {}^{g_1 g_2} \beta^{-1} \cdot b(g_1 g_2)^{-1} \cdot \beta \right) \\
	& = \beta^{-1} \cdot \underbracket{b(g_1) \cdot {}^{g_1} b(g_2) \cdot b(g_1 g_2)^{-1}}_{= a(g_1, g_2)} \cdot \beta \\
	& = a(g_1, g_2).
\end{align*}

% segment-id: o014.li2.chapter6.nonabelian-cohomology.p018
Thus the class of $a$ in $\Hm^2(G,A)$ is determined by
$[c]\in\Hm^1(G,C)$. If $c$ is the constant map with value $1$, then $b$
may also be chosen to be constantly $1$, and likewise for $a$. We have thus
obtained a morphism in $\cate{Set}_\bullet$,
$\delta^1:\Hm^1(G,C)\to\Hm^2(G,A)$.

% segment-id: o014.li2.chapter6.nonabelian-cohomology.th001
\begin{theorem}\label{prop:nonabelian-long}
	Let $1\to A\xrightarrow{u}B\xrightarrow{v}C\to1$ be a short exact
	sequence in $G\dcate{Set}_\bullet$, where $A$ and $B$ are $G$-groups and
	$u$ is a group homomorphism.
% segment-id: o014.li2.chapter6.nonabelian-cohomology.l004
	\begin{enumerate}[(i)]
% segment-id: o014.li2.chapter6.nonabelian-cohomology.i007
		\item If condition \eqref{eqn:pointed-ses-condition} holds, there is
		an exact sequence in $\cate{Set}_\bullet$
% segment-id: o014.li2.chapter6.nonabelian-cohomology.g001
		\begin{equation*}\begin{tikzcd}[row sep=large]
			1 \arrow[r] & \Hm^0(G, A) \arrow[r, "{\Hm^0(u)}"] & \Hm^0(G, B) \arrow[r, "{\Hm^0(v)}"] \arrow[phantom, d, ""{coordinate, name=A}] & \Hm^0(G, C) \arrow[dll, rounded corners, "{\delta^0}" description, to path={
				-- ([xshift=2ex]\tikztostart.east)
				|- (A) [near end]\tikztonodes
				-| ([xshift=-2ex]\tikztotarget.west)
				-- (\tikztotarget)}] \\
			& \Hm^1(G, A) \arrow[r, "{\Hm^1(u)}"] & \Hm^1(G, B). &
		\end{tikzcd}\end{equation*}
% segment-id: o014.li2.chapter6.nonabelian-cohomology.i008
		\item In addition to the assumptions above, suppose that $C$ is also
		a group and $v$ is a group homomorphism; this automatically implies
		\eqref{eqn:pointed-ses-condition}. The exact sequence then extends to
% segment-id: o014.li2.chapter6.nonabelian-cohomology.g002
		\begin{equation*}\begin{tikzcd}[row sep=large]
			1 \arrow[r] & \Hm^0(G, A) \arrow[r, "{\Hm^0(u)}"] & \Hm^0(G, B) \arrow[r, "{\Hm^0(v)}"] \arrow[phantom, d, ""{coordinate, name=A}] & \Hm^0(G, C) \arrow[dll, rounded corners, "{\delta^0}" description, to path={
				-- ([xshift=2ex]\tikztostart.east)
				|- (A) [near end]\tikztonodes
				-| ([xshift=-2ex]\tikztotarget.west)
				-- (\tikztotarget)}] \\
			& \Hm^1(G, A) \arrow[r, "{\Hm^1(u)}"] & \Hm^1(G, B) \arrow[r, "{\Hm^1(v)}"] & \Hm^1(G, C).
		\end{tikzcd}\end{equation*}
% segment-id: o014.li2.chapter6.nonabelian-cohomology.i009
		\item If we further require that $u(A)$ be contained in the center
		$Z_B$ of $B$, then $\delta^0$ is a group homomorphism and the exact
		sequence extends to
% segment-id: o014.li2.chapter6.nonabelian-cohomology.g003
		\begin{equation*}\begin{tikzcd}[row sep=large]
			1 \arrow[r] & \Hm^0(G, A) \arrow[r, "{\Hm^0(u)}"] & \Hm^0(G, B) \arrow[r, "{\Hm^0(v)}"] \arrow[phantom, d, ""{coordinate, name=A}] & \Hm^0(G, C) \arrow[dll, rounded corners, "{\delta^0}" description, to path={
				-- ([xshift=2ex]\tikztostart.east)
				|- (A) [near end]\tikztonodes
				-| ([xshift=-2ex]\tikztotarget.west)
				-- (\tikztotarget)}] \\
			& \Hm^1(G, A) \arrow[r, "{\Hm^1(u)}"] & \Hm^1(G, B) \arrow[r, "{\Hm^1(v)}"] \arrow[phantom, d, ""{coordinate, name=B}] & \Hm^1(G, C) \arrow[dll, rounded corners, "{\delta^1}" description, to path={
				-- ([xshift=2ex]\tikztostart.east)
				|- (B) [near end]\tikztonodes
				-| ([xshift=-2ex]\tikztotarget.west)
				-- (\tikztotarget)}] \\
			& \Hm^2(G, A) . & {} & {}
		\end{tikzcd}\end{equation*}
	\end{enumerate}

	The connecting morphisms $\delta^0$ and $\delta^1$ are defined as above.
	Both are canonical, meaning that they are compatible with morphisms of
	short exact sequences.
\end{theorem}
% segment-id: o014.li2.chapter6.nonabelian-cohomology.pf001
\begin{proof}
	The detailed constructions of $\delta^0$ and $\delta^1$ above already
	show that both are canonical; in case (iii), they also showed that
	$\delta^0$ is a group homomorphism. We now check exactness. Henceforth
	regard $A$ as a subgroup of $B$ and omit $u$.

	In case (i), exactness at $\Hm^0(G,A)$ and $\Hm^0(G,B)$ is clear. At
	$\Hm^0(G,C)$, let $c\in C^G$. By definition, $c$ maps to the basepoint
	of $\Hm^1(G,A)$ if and only if there are $\alpha\in A$ and $b\in B$ such
	that
% segment-id: o014.li2.chapter6.nonabelian-cohomology.q019
	\[ v(b) = c, \quad \forall g \in G,\; b^{-1} \cdot {}^g b = \alpha^{-1} \cdot {}^g \alpha. \]
	But this is equivalent to the existence of $\alpha\in A$ and $b\in B$
	such that $v(b)=c$ and $b\alpha^{-1}\in B^G$, which is in turn
	equivalent to $c\in\Image(\Hm^0(u))$.

	At $\Hm^1(G,A)$, let $a\in Z^1(G,A)$. The class
	$[a]\in\Hm^1(G,A)$ maps to the basepoint if and only if there is a
	$b\in B$ such that $b^{-1}\cdot a(g)\cdot{}^gb=1$ for every $g$. If this
	holds, set $c:=v(b^{-1})$. Applying $v$ to
	${}^gb^{-1}=b^{-1}\cdot a(g)$ shows that $c\in C^G$ and
	$[a]=\delta^0([c])$. The converse argument is similar.

	In case (ii), what remains is exactness at $\Hm^1(G,B)$. Let
	$b\in Z^1(G,B)$. Its class $[b]$ maps to the basepoint if and only if
	there is a $\beta\in B$ such that
	$\beta^{-1}\cdot b(g)\cdot{}^g\beta\in A$ for every $g$. Clearly, this
	is the same as saying that, up to $\sim$, the class comes from
	$Z^1(G,A)$.

	Consider case (iii). Let $c\in Z^1(G,C)$. Its class $[c]$ maps to the
	basepoint if and only if there are maps $b:G\to B$ and $\alpha:G\to A$
	such that $v\circ b=c$ and
% segment-id: o014.li2.chapter6.nonabelian-cohomology.q020
	\[ b(g_1) \cdot {}^{g_1} b(g_2) = \alpha(g_1) \cdot {}^{g_1} \alpha(g_2) \cdot \alpha(g_1 g_2)^{-1} \cdot  b(g_1 g_2). \]
	Since $A\subset Z_B$, if $\alpha^{-1}b:G\to B$ is defined by pointwise
	multiplication, the equation above is also equivalent to
% segment-id: o014.li2.chapter6.nonabelian-cohomology.q021
	\[ (\alpha^{-1} b)(g_1) \cdot {}^{g_1} (\alpha^{-1} b)(g_2) = (\alpha^{-1} b)(g_1 g_2), \]
	that is, $\alpha^{-1}b\in Z^1(G,B)$. Clearly this is equivalent to saying
	that $[c]$ comes from $\Hm^1(G,B)$. This proves exactness.
\end{proof}

% segment-id: o014.li2.chapter6.nonabelian-cohomology.p019
If $A$ and $B$ are further required to be abelian, then
$C\simeq B/u(A)$ is an abelian group, and we recover the familiar version
with coefficients in an abelian group.

% segment-id: o014.li2.chapter6.nonabelian-cohomology.p020
Note that the exact sequence in Theorem~\ref{prop:nonabelian-long} (i)
describes the inverse image of the basepoint under $\Hm^1(u)$, but in
general it cannot determine whether two elements of $\Hm^1(G,A)$ have the
same image under $\Hm^1(u)$. The ``twisting'' construction introduced in
the exercises for this chapter helps address this problem.

% segment-id: o014.li2.chapter6.nonabelian-cohomology.e001
\begin{example}
	Let $B$ be a $G$-group and $A$ a subgroup stable under the $G$-action.
	Equip the set $B/A$ with the basepoint $1\cdot A$ and a left $G$-action.
	This gives a short exact sequence in $G\dcate{Set}_\bullet$
% segment-id: o014.li2.chapter6.nonabelian-cohomology.q022
	\[ 1 \to A \xrightarrow{\text{inclusion}} B \xrightarrow{\text{quotient}} B/A \to 1. \]
	Condition \eqref{eqn:pointed-ses-condition} clearly holds. Hence
	Theorem~\ref{prop:nonabelian-long} (i) gives an exact sequence in
	$\cate{Set}_\bullet$
% segment-id: o014.li2.chapter6.nonabelian-cohomology.q023
	\[ 1 \to A^G \to B^G \to (B/A)^G \xrightarrow{\delta^0} \Hm^1(G, A) \to \Hm^1(G, B). \]

	The evident map $B^G/A^G\hookrightarrow(B/A)^G$ is injective, but it is
	often not bijective: a coset stable under the $G$-action need not contain
	a $G$-fixed point; the exercises will give an example. Nevertheless, the
	exact sequence identifies the cohomological obstruction. For a morphism
	$f$ in $\cate{Set}_\bullet$, define $\Ker(f)$ to be the inverse image of
	the basepoint. Then
% segment-id: o014.li2.chapter6.nonabelian-cohomology.q024
	\begin{align*}
		% Editorial correction: the frozen Chinese source and Indonesian witness
		% switch to \delta_0 here; the connecting morphism defined above is \delta^0.
		\text{the coset}\; bA \in (B/A)^G \;\text{comes from}\; B^G & \iff bA \in \Ker(\delta^0), \\
		(B/A)^G = B^G/A^G & \iff \Ker[\Hm^1(G, A) \to \Hm^1(G, B)] = 1.
	\end{align*}
\end{example}

% segment-id: o014.li2.chapter6.nonabelian-cohomology.p021
We close this section with the nonabelian version of Shapiro's Lemma
(Theorem~\ref{prop:Shapiro}). As preparation, first define a nonabelian
version of induced modules using a mapping space; compare
Lemma~\ref{prop:Ind-as-mappings}. The construction applies to objects of
$H\dcate{Grp}$ or $H\dcate{Set}_\bullet$, where $H$ is a subgroup of $G$;
the group action is written $(h,a)\mapsto{}^ha$.

% segment-id: o014.li2.chapter6.nonabelian-cohomology.d005
\begin{definition}
	Let $H$ be a subgroup of $G$ and let $A$ be an object of
	$H\dcate{Grp}$ (or $H\dcate{Set}_\bullet$). Define the object
	$\Ind^G_H(A)$ of $G\dcate{Grp}$ (or $G\dcate{Set}_\bullet$) as follows.
	As an object of $G\dcate{Set}$, it is
% segment-id: o014.li2.chapter6.nonabelian-cohomology.q025
	\begin{gather*}
		\Ind^G_H(A) := \left\{\begin{array}{r|l}
			f: G \to A & \forall h \in H, \; \forall x \in G \\
			& f(hx) = {}^h f(x).
		\end{array}\right\}, \\
		({}^g f)(x) := f(xg), \quad f \in \Ind^G_H(A), \; x, g \in G,
	\end{gather*}
	with the group operation (respectively, basepoint) defined pointwise on
	maps (respectively, by the constant map to the basepoint). We call
	$A\mapsto\Ind^G_H(A)$ induction.
\end{definition}

% segment-id: o014.li2.chapter6.nonabelian-cohomology.p022
The induction construction gives compatible functors
$H\dcate{Grp}\to G\dcate{Grp}$ and
$H\dcate{Set}_\bullet\to G\dcate{Set}_\bullet$.

% segment-id: o014.li2.chapter6.nonabelian-cohomology.th002
\begin{theorem}\label{prop:Shapiro-noncommutative}
	\index{Shapiro's lemma}
	Let $H$ be a subgroup of $G$ and let $A$ be an object of
	$H\dcate{Set}_\bullet$ or $H\dcate{Grp}$. There are canonical
	isomorphisms in $\cate{Set}_\bullet$
% segment-id: o014.li2.chapter6.nonabelian-cohomology.q026
	\[ \Hm^i\left(G, \Ind^G_H(A)\right) \rightiso \Hm^i(H, A), \]
	with $i=0$ if $A$ comes from $H\dcate{Set}_\bullet$, or $i=0,1$ if $A$
	comes from $H\dcate{Grp}$. The maps are as follows.
% segment-id: o014.li2.chapter6.nonabelian-cohomology.l005
	\begin{description}
% segment-id: o014.li2.chapter6.nonabelian-cohomology.i010
		\item[($i=0$)] Take the map $f\mapsto f(1_G)$, which defines
		$Z^0\left(G,\Ind^G_H(A)\right)=\Ind^G_H(A)^G\to A^H=Z^0(H,A)$.
% segment-id: o014.li2.chapter6.nonabelian-cohomology.i011
		\item[($i=1$)] Take the map
		$Z^1\left(G,\Ind^G_H(A)\right)\to Z^1(H,A)$ sending
		$c:G\to\Ind^G_H(A)$ to
% segment-id: o014.li2.chapter6.nonabelian-cohomology.q027
		\begin{align*}
			(c|_H)(\cdot)(1_G): H & \to A \\
			h & \mapsto c(h)(1_G).
		\end{align*}
	\end{description}
	Both induce morphisms on $\Hm^i$. If $A$ comes from $H\dcate{Grp}$, the
	isomorphism for $i=0$ is also an isomorphism of groups.
\end{theorem}
% segment-id: o014.li2.chapter6.nonabelian-cohomology.pf002
\begin{proof}
	Routine arguments show that the maps in both cases are well-defined and
	preserve the basepoint or group structure; details are left to the reader.

	For $i=0$, it is easy to see that $\Ind^G_H(A)^G\to A^H$ is indeed a
	bijection and gives an isomorphism of groups if $A$ is an $H$-group. We
	now focus on the case where $A$ is an object of $H\dcate{Grp}$ and
	$i=1$.

	Suppose that $r,s\in Z^1(G,\Ind^G_H(A))$ have the same image in
	$\Hm^1(H,A)$. We prove that $r\sim s$. This assumption is equivalent to
	the existence of an $\alpha\in A$ such that
% segment-id: o014.li2.chapter6.nonabelian-cohomology.q028
	\[ r(h)(1_G) = \alpha^{-1} \cdot s(h)(1_G) \cdot {}^h \alpha \]
	for every $h\in H$. There is a $\zeta\in\Ind^G_H(A)$ with
	$\zeta(1_G)=\alpha$. Adjusting $s$ by $\zeta$ within its equivalence
	class, we may ensure that $r(h)(1_G)=s(h)(1_G)$ for every $h$.

	The definition of $Z^1(G,\Ind^G_H(A))$ implies the identities in $A$
% segment-id: o014.li2.chapter6.nonabelian-cohomology.q029
	\begin{gather*}
		r(g_1 g_2)(x) =r(g_1)(x) \cdot r(g_2)(xg_1), \quad s(g_1 g_2)(x) = s(g_1)(x) \cdot s(g_2)(xg_1)
	\end{gather*}
	for all $x,g_1,g_2\in G$. On the one hand, substituting $g_1=g\in G$
	and $g_2=g^{-1}x^{-1}$ and rearranging gives
% segment-id: o014.li2.chapter6.nonabelian-cohomology.q030
	\begin{equation}\label{eqn:Shapiro-noncommutative-aux0}
		\begin{aligned}
			r(g)(x) & = r(x^{-1})(x) \cdot r(g^{-1} x^{-1})(xg)^{-1}, \\
			s(g)(x) & = s(x^{-1})(x) \cdot s(g^{-1} x^{-1})(xg)^{-1}.
		\end{aligned}
	\end{equation}
	On the other hand, substituting $g_1=x^{-1}$ and $g_2=h\in H$ gives
% segment-id: o014.li2.chapter6.nonabelian-cohomology.q031
	\begin{equation}\label{eqn:Shapiro-noncommutative-aux1}
		s(x^{-1} h)(x) \cdot r(x^{-1} h)(x)^{-1} = s(x^{-1})(x) \cdot r(x^{-1})(x)^{-1}.
	\end{equation}

	For every $x\in G$, define
	$t(x):=s(x^{-1})(x)\cdot r(x^{-1})(x)^{-1}$. We claim that
	$t:G\to A$ belongs to $\Ind^G_H(A)$. Indeed, if $h\in H$ and $x\in G$,
	then
% segment-id: o014.li2.chapter6.nonabelian-cohomology.q032
	\begin{multline*}
		t(hx) = s(x^{-1} h^{-1})(hx) \cdot r(x^{-1} h^{-1})(hx)^{-1} \\
		= {}^h \left( s(x^{-1}h^{-1})(x) \cdot r(x^{-1} h^{-1})(x)^{-1} \right) \\
		\xlongequal{\text{\eqref{eqn:Shapiro-noncommutative-aux1}}} {}^h \left( s(x^{-1})(x) \cdot r(x^{-1})(x)^{-1} \right) = {}^h t(x).
	\end{multline*}

	To prove $r\sim s$, it is enough to check the identity in
	$\Ind^G_H(A)$ given by $t^{-1}\cdot s(g)\cdot{}^gt=r(g)$. The value of
	the left-hand side at $x\in G$ is
% segment-id: o014.li2.chapter6.nonabelian-cohomology.q033
	\[ r(x^{-1})(x) \cdot s(x^{-1})(x)^{-1} \cdot s(g)(x) \cdot s(g^{-1} x^{-1})(xg) \cdot r(g^{-1} x^{-1})(xg)^{-1}; \]
	expanding $s(g)(x)$ and $r(g)(x)$ using
	\eqref{eqn:Shapiro-noncommutative-aux0}, this expression is visibly equal
	to $r(g)(x)$. This proves injectivity.

	Finally, prove surjectivity. Let $s^\flat\in Z^1(H,A)$. Write $\pi$ for
	the quotient map $G\twoheadrightarrow H\backslash G$, and choose any map
	$\sigma:H\backslash G\to G$ satisfying $\pi\sigma=\identity$. Without
	loss of generality, assume $\sigma(H\cdot1_G)=1_G$. Every $g\in G$ has a
	unique representation
% segment-id: o014.li2.chapter6.nonabelian-cohomology.q034
	\[ g = \tau(g) \sigma(Hg), \quad \tau(g) \in H. \]

	Note that $\tau(1_G)=1_G$. For every $g,x\in G$, define
% segment-id: o014.li2.chapter6.nonabelian-cohomology.q035
	\[ s(g)(x) := {}^{\tau(x)} s^\flat\left( \tau(\sigma(Hx)g) \right) \; \in A. \]
	If $h\in H$, then $\tau(hx)=h\tau(x)$ immediately gives
	$s(g)(hx)={}^hs(g)(x)$, so $s(g)\in\Ind^G_H(A)$. The next step is to
	show that $s\in Z^1(G,\Ind^G_H(A))$. The reader is first asked to verify
% segment-id: o014.li2.chapter6.nonabelian-cohomology.q036
	\begin{equation}\label{eqn:Shapiro-noncommutative-aux2}
		\tau(uv) = \tau(u) \tau(\sigma(Hu)v), \quad u, v \in G.
	\end{equation}

	We need to prove that
	$s(g_1g_2)(x)=s(g_1)(x)\cdot s(g_2)(xg_1)$ for all
	$g_1,g_2,x\in G$. By the definition of $Z^1(H,A)$, the left-hand side
	equals
% segment-id: o014.li2.chapter6.nonabelian-cohomology.q037
	\begin{multline*}
		{}^{\tau(x)} s^\flat\left( \tau(\sigma(Hx)g_1 g_2) \right) \xlongequal{\text{\eqref{eqn:Shapiro-noncommutative-aux2}}} {}^{\tau(x)} s^\flat\left( \tau(\sigma(Hx)g_1) \tau(\sigma(Hx g_1)g_2) \right) \\
		= {}^{\tau(x)} s^\flat\left( \tau(\sigma(Hx)g_1) \right) \cdot {}^{\tau(x)\tau(\sigma(Hx)g_1)} s^\flat\left( \tau(\sigma(Hxg_1)g_2) \right) \\
		\xlongequal{\text{\eqref{eqn:Shapiro-noncommutative-aux2}}} {}^{\tau(x)} s^\flat\left( \tau(\sigma(Hx)g_1) \right) \cdot {}^{\tau(xg_1)} s^\flat\left( \tau(\sigma(Hxg_1)g_2) \right),
	\end{multline*}
	which is the right-hand side.

	Moreover, if $h\in H$, then $s(h)(1_G)=s^\flat(h)$, so $s$ maps to
	$s^\flat$. This proves surjectivity.
\end{proof}

% segment-id: o014.li2.chapter6.nonabelian-cohomology.p023
Clearly the isomorphism in Theorem~\ref{prop:Shapiro-noncommutative}, when
% Editorial correction: the theorem starts with an H-group; the frozen source
% and Indonesian witness call A a G-group in this concluding sentence.
$A$ is an abelian $H$-group, reduces to the special case of
Proposition~\ref{prop:Shapiro-std-coh}. Moreover, a direct computation shows
that these isomorphisms are compatible with the connecting morphisms in
Theorem~\ref{prop:nonabelian-long}; the details are left as an exercise.

% segment-id: o014.li2.chapter6.nonabelian-cohomology.r002
\begin{remark}[Nonabelian cohomology of profinite groups]\label{rem:pro-finite-coh-nonabelian}
	\index{$G$-group!smooth}
	Let $G$ be a profinite group. An object $A$ of
	$G\dcate{Set}_\bullet$ or $G\dcate{Grp}$ is called \emph{smooth} if
	$A=\bigcup_KA^K$, where $K$ ranges over the open normal subgroups of $G$;
	in this case $A$ is given the discrete topology. Since all arguments in
	this section are based on analogies with the standard cochain complex, all
	the preceding results apply to a profinite group $G$ and smooth $A$. The
	only difference is that $c\in Z^i(G,A)$, as a map from $G$ to $A$, must
	be continuous, that is, locally constant. Thus $\Hm^i(G,A)$ and the
	various canonical morphisms can be defined for $i=0,1$, and
% segment-id: o014.li2.chapter6.nonabelian-cohomology.q038
	\begin{align*}
		Z^i(G, A) & = \varinjlim_K Z^i(G/K, A^K), \\
		\Hm^i(G, A) & = \varinjlim_K \Hm^i(G/K, A^K), \quad i = 0, 1.
	\end{align*}

	In addition, the definition of $\Ind^G_H(A)$ must require $f:G\to A$ to
	be locally constant. The proof of
	Theorem~\ref{prop:Shapiro-noncommutative} also requires
	Lemma~\ref{prop:profinite-section} to ensure that the $\zeta:G\to A$ in
	the injectivity part can be chosen locally constant and that the
	$\sigma:H\backslash G\to G$ in the surjectivity part is continuous.

	Finally, Remark~\ref{rem:Z1-vs-section} remains valid for profinite groups
	$G$: $Z^1(G,A)$ corresponds to the set of continuous sections of the
	projection homomorphism $\pi:A\rtimes G\to G$, while $\Hm^1(G,A)$ is the
	quotient of the set of continuous sections by the conjugation action of
	$A$.
\end{remark}
